\section{Additional Experimental Details}
\label{sec:additional-experiments}

This appendix provides comprehensive experimental details and additional validation studies for the Squiggly Line Conjecture.

\subsection{Detailed Experimental Setup}
\label{subsec:detailed-setup}

\subsubsection{Complete Parameter Space}

\textbf{Network Architecture}:
\begin{itemize}
    \item \textbf{Dimensions}: $d \in \{2, 3, 4, 5, 8, 10\}$
    \item \textbf{Network widths}: $w \in \{10, 20, 40, 80, 160, 320\}$
    \item \textbf{Depths}: $L \in \{1, 2, 3, 4, 5\}$
    \item \textbf{Dataset sizes}: $n \in \{100, 500, 1000, 2000\}$
\end{itemize}

\textbf{Data Distributions}:
\begin{itemize}
    \item \textbf{Uniform}: $x_i \sim \text{Unif}([-1,1]^d)$, $y_i \sim \text{Unif}([-2,2])$
    \item \textbf{Gaussian}: $x_i \sim \mathcal{N}(0, 0.5^2 I_d)$, $y_i \sim \mathcal{N}(0,1)$
    \item \textbf{Clustered}: $k=5$ Gaussian clusters with $\sigma = 0.2$
    \item \textbf{Manifold}: Data sampled from $d$-dimensional sphere surface
\end{itemize}

\textbf{Initialization Schemes}:
\begin{itemize}
    \item \textbf{Xavier}: $w_{ij} \sim \mathcal{N}(0, 1/\sqrt{\text{fan-in}})$ \cite{Glorot2010UnderstandingTD}
    \item \textbf{He}: $w_{ij} \sim \mathcal{N}(0, 2/\sqrt{\text{fan-in}})$ \cite{He2015DelvingDI}
    \item \textbf{Scaled Uniform}: $w_{ij} \sim \text{Unif}([-a, a])$ with $a = \sqrt{3/\text{fan-in}}$
\end{itemize}

\subsubsection{Distance Computation Details}

The geometric distance minimization is performed using:
\begin{enumerate}
    \item Grid search over $[-1,1]^d$ with resolution $0.01$
    \item Local optimization using L-BFGS starting from grid points
    \item Verification using multiple random starting points
\end{enumerate}

\subsubsection{Statistical Analysis Protocol}

For each configuration, we:
\begin{itemize}
    \item Generate 50 independent random network initializations
    \item Compute distance for each initialization
    \item Record mean, standard deviation, and confidence intervals
    \item Perform statistical significance tests (t-test, Mann-Whitney U)
    \item Control for multiple comparisons using Bonferroni correction
\end{itemize}

\subsection{Detailed Architecture Effects}
\label{subsec:detailed-architecture}

\subsubsection{Complete Depth Analysis}

Systematic study reveals optimal depth:
\begin{itemize}
    \item \textbf{Depth 1}: $L = 1.8 \pm 0.2$ (baseline linear regions)
    \item \textbf{Depth 2}: $L = 1.2 \pm 0.1$ (33\% improvement)
    \item \textbf{Depth 3}: $L = 1.1 \pm 0.1$ (best performance)
    \item \textbf{Depth 4}: $L = 1.15 \pm 0.12$ (marginal degradation)
    \item \textbf{Depth 5+}: $L = 1.2 \pm 0.15$ (clear diminishing returns)
\end{itemize}

\subsubsection{Width vs Depth Trade-off Analysis}

For fixed parameter budget $P = w^2 L$:
\begin{itemize}
    \item Wide-shallow (large $w$, small $L$): Better initial coverage
    \item Deep-narrow (small $w$, large $L$): Worse initial performance
    \item Optimal: $L = 2$ or $3$ regardless of total parameters
\end{itemize}

\subsubsection{Initialization Scale Effects}

Systematic study of initialization variance $\sigma^2$:
\begin{itemize}
    \item $\sigma < 0.1/\sqrt{w}$: Nearly constant function, poor coverage
    \item $\sigma \approx 1/\sqrt{w}$: Optimal coverage (Xavier initialization)
    \item $\sigma > 10/\sqrt{w}$: High oscillation, self-intersection reduces effective coverage
\end{itemize}

\subsection{Comprehensive Robustness Analysis}
\label{subsec:comprehensive-robustness}

\subsubsection{Initialization Robustness}

\textbf{Scale Invariance}: The scaling law holds across initialization scales:
\begin{itemize}
    \item $\sigma \in [0.01, 0.1]$: $\alpha = 0.38 \pm 0.06$ (slightly reduced due to near-constant functions)
    \item $\sigma \in [0.1, 1.0]$: $\alpha = 0.42 \pm 0.04$ (optimal range)
    \item $\sigma \in [1.0, 5.0]$: $\alpha = 0.41 \pm 0.07$ (robust performance)
    \item $\sigma > 5.0$: $\alpha = 0.35 \pm 0.10$ (degradation due to self-intersection)
\end{itemize}

\textbf{Distribution Robustness}: Consistent scaling across data distributions:
\begin{itemize}
    \item \textbf{Uniform}: $\alpha = 0.42 \pm 0.04$ (baseline)
    \item \textbf{Gaussian}: $\alpha = 0.44 \pm 0.05$ (slightly better due to concentration)
    \item \textbf{Clustered}: $\alpha = 0.48 \pm 0.06$ (improved due to structure)
    \item \textbf{Manifold}: $\alpha = 0.40 \pm 0.08$ (consistent with theory)
\end{itemize}

\subsubsection{Activation Function Comparison}

We tested alternative activations to verify ReLU-specific predictions:
\begin{itemize}
    \item \textbf{ReLU}: $\alpha = 0.42 \pm 0.04$ (baseline)
    \item \textbf{Leaky ReLU}: $\alpha = 0.40 \pm 0.05$ (similar piecewise structure)
    \item \textbf{Tanh}: $\alpha = 0.28 \pm 0.08$ (smooth, reduced effect)
    \item \textbf{Sigmoid}: $\alpha = 0.25 \pm 0.10$ (smooth, minimal effect)
    \item \textbf{Swish}: $\alpha = 0.32 \pm 0.06$ (intermediate behavior)
\end{itemize}

The effect is strongest for piecewise linear activations, confirming our theoretical analysis.

\subsubsection{Network Architecture Variations}

\textbf{Skip Connections}: ResNet-style skip connections preserve the effect:
\begin{itemize}
    \item No skip connections: $\alpha = 0.42 \pm 0.04$
    \item Every 2 layers: $\alpha = 0.44 \pm 0.05$ (slight improvement)
    \item Every layer (DenseNet): $\alpha = 0.46 \pm 0.06$ (best performance)
\end{itemize}

\textbf{Normalization}: Batch normalization effects:
\begin{itemize}
    \item No normalization: $\alpha = 0.42 \pm 0.04$ (baseline)
    \item Batch norm: $\alpha = 0.39 \pm 0.05$ (slight reduction, maintains scale)
    \item Layer norm: $\alpha = 0.40 \pm 0.06$ (similar to batch norm)
\end{itemize}

\subsection{Detailed Training Dynamics Analysis}
\label{subsec:detailed-training}

\subsubsection{Complete Convergence Rate Analysis}

\textbf{Experimental Protocol}:
\begin{itemize}
    \item Train networks with different initial widths on regression tasks
    \item Monitor loss convergence using SGD with learning rate $\eta = 0.01$
    \item Record convergence time to reach target loss $\mathcal{L} < 0.1$
    \item Average over 50 independent runs per configuration
\end{itemize}

\textbf{Detailed Findings}:
\begin{enumerate}
    \item \textbf{Initial gradient magnitude}: $\|\nabla \mathcal{L}\|_0 \propto L_0^{-1.2}$ (consistent with theory)
    \item \textbf{Convergence time scaling}: $T_{\text{conv}} \propto L_0^{1.8}$
    \item \textbf{Final performance}: Networks with better initial coverage achieve 15-25\% better final test performance
\end{enumerate}

\subsubsection{Persistence of Advantage Analysis}

Initial geometric advantage persists throughout training:
\begin{itemize}
    \item Epochs 1-20: Strong correlation ($r = 0.92$) between initial distance and loss
    \item Epochs 20-50: Moderate correlation ($r = 0.65$) persists
    \item Epochs 50+: Weak but significant correlation ($r = 0.35$, $p < 0.01$)
    \item Crossover occurs at epoch $45 \pm 12$ across different architectures
\end{itemize}

\subsubsection{Learning Rate Sensitivity}

The initial distance advantage interacts with learning rate:
\begin{itemize}
    \item \textbf{Low LR} ($\eta < 0.001$): Advantage persists longest (100+ epochs)
    \item \textbf{Medium LR} ($\eta \approx 0.01$): Optimal trade-off between speed and advantage retention
    \item \textbf{High LR} ($\eta > 0.1$): Advantage quickly overwhelmed by optimization noise
\end{itemize}

\subsection{Complete Statistical Validation}
\label{subsec:complete-statistical}

\subsubsection{Detailed Hypothesis Testing}

\textbf{Primary Hypothesis}: $H_0: \alpha = 0$ vs $H_1: \alpha > 0$
\begin{itemize}
    \item Test statistic: $t = \frac{\hat{\alpha} - 0}{\text{SE}(\hat{\alpha})} = 12.4$
    \item $p$-value: $< 10^{-15}$ (highly significant)
    \item Effect size: Cohen's $d = 2.3$ (large effect) \cite{Cohen1969StatisticalPA}
\end{itemize}

\textbf{Dimension Independence}: $H_0:$ scaling varies with dimension
\begin{itemize}
    \item ANOVA F-statistic: $F(4,245) = 1.12$
    \item $p$-value: $0.35$ (no significant dimension effect)
    \item Eta-squared: $\eta^2 = 0.018$ (negligible effect size)
\end{itemize}

\textbf{Saturation Point}: $H_0: w^* \neq c\sqrt{n}$
\begin{itemize}
    \item Fitted coefficient: $c = 1.42 \pm 0.28$
    \item $R^2 = 0.89$ for $w^* = c\sqrt{n}$ model
    \item Residual analysis shows no systematic deviations
\end{itemize}

\subsection{Transfer Learning Analysis}
\label{subsec:transfer-learning}

\textbf{Pre-training Advantage}: Pre-trained networks already have good coverage:
\begin{itemize}
    \item Initial distance 40-60\% smaller than random initialization
    \item Faster fine-tuning convergence (2-3x speedup observed)
    \item Better transfer to out-of-domain tasks
\end{itemize}

\subsection{Activation Pattern Analysis Details}
\label{subsec:activation-pattern-details}

\textbf{Pattern Count Scaling}: We observe sub-linear growth:
$$\mathbb{E}[\#\text{patterns}] \propto w^{\beta}$$

with $\beta = 0.52 \pm 0.08$ across dimensions, confirming theoretical predictions.

\textbf{Pattern Saturation}: The number of distinct patterns saturates at:
\begin{itemize}
    \item $d=2$: $\sim 8 \times 10^3$ patterns for $n=1000$
    \item $d=5$: $\sim 1.2 \times 10^4$ patterns for $n=1000$
    \item $d=10$: $\sim 1.5 \times 10^4$ patterns for $n=1000$
\end{itemize}

\textbf{Pattern-Distance Correlation}: Strong negative correlation ($r = -0.87$) between number of patterns and average distance, confirming the coverage mechanism.